\documentclass[11pt]{article}

% ============================================================================
% PACKAGES
% ============================================================================
\usepackage[utf8]{inputenc}
\usepackage[T1]{fontenc}
\usepackage{amsmath,amssymb,amsthm}
\usepackage{algorithm}
\usepackage{algpseudocode}
\usepackage{booktabs}
\usepackage{graphicx}
\usepackage{hyperref}
\usepackage{xcolor}
\usepackage[margin=1in]{geometry}
\usepackage{caption}
\usepackage{subcaption}
\usepackage{enumitem}
\usepackage{natbib}

% ============================================================================
% THEOREM ENVIRONMENTS
% ============================================================================
\newtheorem{theorem}{Theorem}
\newtheorem{lemma}[theorem]{Lemma}
\newtheorem{proposition}[theorem]{Proposition}
\newtheorem{corollary}[theorem]{Corollary}
\newtheorem{definition}[theorem]{Definition}
\newtheorem{remark}{Remark}

% ============================================================================
% TITLE
% ============================================================================
\title{Completion-Aware Decreasing: A Simple Polynomial-Time Heuristic\\for Bin Packing with Reduced Waste}

\author{
  Kathan Sanghavi\thanks{Corresponding author. Email: \texttt{kathanchess1396@gmail.com}}
}

\date{March 2026}

\begin{document}
\maketitle

% ============================================================================
% ABSTRACT
% ============================================================================
\begin{abstract}
We introduce \textbf{Completion-Aware Decreasing (CAD)}, a simple $O(n^2)$-time heuristic
for the one-dimensional bin packing problem. CAD augments Best-Fit Decreasing (BFD)
with an immediate post-placement \emph{completion step}: after assigning each item to a bin,
it searches the remaining unplaced items for one or two items that tightly fill the bin's
residual capacity. This breaks the strict size-ordered processing of classical decreasing
heuristics, enabling better item pairings that reduce total waste. In experiments on
1{,}400+ instances across six instance classes, CAD uses fewer bins than First-Fit
Decreasing (FFD) on every hard instance tested (162 wins, 3 losses on hard instances
with $n \geq 200$, $p < 10^{-33}$, Wilcoxon signed-rank test) and achieves a 55\%
reduction in waste relative to the $L_2$ lower bound on Scholl-type hard instances.
CAD also outperforms Modified FFD (MFFD) on all hard and medium instance types.
The algorithm requires no parameter tuning for practical use and runs in under 400ms
for $n = 1{,}000$.
\end{abstract}

\noindent\textbf{Keywords:} bin packing, heuristic, combinatorial optimization, first-fit decreasing, approximation algorithm

% ============================================================================
% 1. INTRODUCTION
% ============================================================================
\section{Introduction}

The one-dimensional bin packing problem (BPP) is among the most studied problems in
combinatorial optimization. Given $n$ items with sizes $s_1, \ldots, s_n \in (0, C]$
and bins of capacity $C$, the objective is to pack all items into the minimum number
of bins. The problem is NP-hard in the strong sense \citep{garey1979computers}, and
consequently decades of work have focused on polynomial-time approximation heuristics.

The most widely used heuristic family is the \emph{decreasing} family, which sorts
items by non-increasing size and then applies a placement rule:
\begin{itemize}[nosep]
  \item \textbf{First-Fit Decreasing (FFD):} Place each item in the first bin with
        sufficient remaining capacity. Asymptotic ratio $\frac{11}{9}\text{OPT} + \frac{6}{9}$
        \citep{dosa2007tight}.
  \item \textbf{Best-Fit Decreasing (BFD):} Place each item in the bin with the
        \emph{tightest} remaining capacity. Same asymptotic bound as FFD.
  \item \textbf{Modified FFD (MFFD):} Classifies items into four size categories and
        applies specialized placement rules. Achieves $\frac{71}{60}\text{OPT} + 1$
        \citep{johnson1985bin}.
\end{itemize}

All of these heuristics process items in strict decreasing-size order. Once an item is
placed, the algorithm moves to the next item without considering whether the bin just
used could be immediately completed with items from later in the sequence. This is
a missed opportunity: when a large item leaves significant residual capacity, a good
completing item may exist among the remaining unplaced items, but it will not be
considered until its turn in the decreasing sequence---by which point it may be
consumed by a different bin.

\paragraph{Our contribution.} We propose \textbf{Completion-Aware Decreasing (CAD)},
which augments BFD with an immediate completion search after each placement. The
completion step looks for one or two unplaced items that fill the bin's residual capacity
tightly (within a small waste tolerance). This simple modification yields substantial
empirical improvements on hard instances while preserving the $O(n^2)$ time complexity
of BFD.

The key idea is related to \emph{bin completion} \citep{korf2002new,fukunaga2007bin},
which also considers filling bins to capacity. However, bin completion is an
\emph{exact} algorithm with exponential worst-case time. CAD, by contrast, is a
greedy polynomial-time heuristic that performs a single pass: it never backtracks and
never reopens a bin once it moves on. This makes CAD practical for the same
applications where FFD is currently used.

% ============================================================================
% 2. ALGORITHM
% ============================================================================
\section{The CAD Algorithm}

\subsection{Notation}

Let $I = \{1, \ldots, n\}$ be the set of items with sizes $s_i \in (0, C]$ for bin
capacity $C$. A \emph{bin} $B_j$ is a subset of $I$ with $\sum_{i \in B_j} s_i \leq C$.
The \emph{residual capacity} of bin $B_j$ is $r_j = C - \sum_{i \in B_j} s_i$.

\subsection{Algorithm Description}

\begin{algorithm}[H]
\caption{Completion-Aware Decreasing (CAD)}
\label{alg:cad}
\begin{algorithmic}[1]
\Require Items $s_1, \ldots, s_n$; bin capacity $C$; thresholds $\alpha = 0.15$, $\beta = 0.10$, $\gamma = 0.25$
\Ensure A valid bin packing
\State Sort items by decreasing size; let $\sigma$ be the sorted order
\State $\texttt{used}[i] \gets \texttt{false}$ for all $i$
\For{each $i$ in order $\sigma$}
  \If{$\texttt{used}[i]$} \textbf{continue} \EndIf
  \State $\texttt{used}[i] \gets \texttt{true}$
  \State Place $i$ using Best-Fit into bin $B_j$
  \State $r \gets$ residual capacity of $B_j$
  \If{$r = 0$} \textbf{continue} \EndIf
  \Statex \hspace{\algorithmicindent}\textit{// Phase 1: Single-item completion}
  \State $k^* \gets \arg\min_{k : \neg\texttt{used}[k],\, s_k \leq r} (r - s_k)$
  \If{$r - s_{k^*} \leq \alpha \cdot r$}
    \State Place $k^*$ in $B_j$; $\texttt{used}[k^*] \gets \texttt{true}$; $r \gets r - s_{k^*}$
  \EndIf
  \If{$r = 0$} \textbf{continue} \EndIf
  \Statex \hspace{\algorithmicindent}\textit{// Phase 2: Two-item completion}
  \State Search pairs $(k_1, k_2)$ with $\neg\texttt{used}[k_i]$, $s_{k_1} + s_{k_2} \leq r$, $s_{k_1} \geq \gamma \cdot r$
  \State Let $(k_1^*, k_2^*) = \arg\min (r - s_{k_1} - s_{k_2})$ over feasible pairs
  \If{$r - s_{k_1^*} - s_{k_2^*} \leq \beta \cdot r$}
    \State Place $k_1^*, k_2^*$ in $B_j$; mark both used
  \EndIf
\EndFor
\State Place any remaining unused items via First-Fit
\end{algorithmic}
\end{algorithm}

\subsection{Complexity}

Each of the $n$ items in the main loop triggers at most one scan of remaining items
for single-item completion ($O(n)$) and one scan of pairs for two-item completion.
The pair scan is pruned by requiring $s_{k_1} \geq \gamma \cdot r$ and using the
sorted order to break early, yielding $O(n)$ amortized per item. The overall time
complexity is $O(n^2)$, which is the same as BFD. Space complexity is $O(n)$.

\subsection{Threshold Selection}

The algorithm uses three parameters:
\begin{itemize}[nosep]
  \item $\alpha = 0.15$: maximum waste fraction for single-item completion
  \item $\beta = 0.10$: maximum waste fraction for two-item completion
  \item $\gamma = 0.25$: minimum size ratio for pair search (pruning)
\end{itemize}

These values were chosen conservatively. The single-item threshold $\alpha$ is more
generous than the pair threshold $\beta$ because a single completing item is less
likely to ``steal'' a useful item from a later bin. The pruning parameter $\gamma$
excludes pairs where one item is very small, as such pairs are unlikely to tightly fill
the residual capacity. We show in Section~\ref{sec:experiments} that performance is
robust across a range of threshold values.

% ============================================================================
% 3. EXPERIMENTAL SETUP
% ============================================================================
\section{Experimental Setup}
\label{sec:experiments}

\subsection{Instance Classes}

We evaluate on six instance classes spanning a range of difficulty regimes,
following the conventions of \citet{scholl1997bison} and \citet{falkenauer1996hybrid}:

\begin{table}[h]
\centering
\caption{Instance classes used in experiments.}
\label{tab:instances}
\begin{tabular}{llcl}
\toprule
\textbf{Class} & \textbf{Item Range} & \textbf{Capacity} & \textbf{Characteristics} \\
\midrule
\textsc{Hard-150}    & $[20, 100]$ & 150  & Scholl-type hard instances \\
\textsc{Hard-200}    & $[30, 150]$ & 200  & Larger hard instances \\
\textsc{Medium}      & $[10, 80]$  & 100  & Medium-sized items \\
\textsc{Large-Items} & $[250, 500]$& 1000 & Very large items \\
\textsc{Uniform}     & $[1, 99]$   & 100  & Standard uniform \\
\textsc{Triplet}     & Structured  & 1000 & Falkenauer-style; OPT $= n/3$ \\
\bottomrule
\end{tabular}
\end{table}

\subsection{Baselines}

We compare against three standard heuristics:
\begin{itemize}[nosep]
  \item \textbf{FFD} \citep{dosa2007tight}: $O(n \log n)$, the de facto industry standard.
  \item \textbf{BFD}: $O(n^2)$, tightest-fit variant of FFD.
  \item \textbf{MFFD} \citep{johnson1985bin}: $O(n \log n)$, four-category item classification with the best known polynomial-time asymptotic ratio of $\frac{71}{60}$.
\end{itemize}

\subsection{Methodology}

For each (class, $n$) configuration, we generate 50 independent random instances using
different seeds. We measure:
\begin{itemize}[nosep]
  \item \textbf{Approximation ratio}: $\text{bins}(A) / L_2$, where $L_2 = \lceil \sum s_i / C \rceil$ is the continuous lower bound.
  \item \textbf{Win/tie/loss counts}: number of instances where CAD uses fewer, equal, or more bins than each baseline.
  \item \textbf{Statistical significance}: Wilcoxon signed-rank test across paired instances.
\end{itemize}

Item sizes $n$ range from 60 to 500. All experiments use integer item sizes generated
uniformly at random within the specified ranges.

% ============================================================================
% 4. RESULTS
% ============================================================================
\section{Results}

\subsection{Main Result: Hard Instances}

Table~\ref{tab:hard150} shows CAD vs.\ FFD on \textsc{Hard-150} instances.
CAD achieves zero losses across all 250 instances and the improvement grows with $n$.

\begin{table}[h]
\centering
\caption{CAD vs.\ FFD on \textsc{Hard-150} instances (50 trials per $n$).}
\label{tab:hard150}
\begin{tabular}{rcccccc}
\toprule
$n$ & FFD Ratio & CAD Ratio & Wins & Ties & Losses & $p$-value \\
\midrule
60  & 1.0208 & 1.0085 & 15 & 35 & 0 & $1.08 \times 10^{-4}$ \\
100 & 1.0198 & 1.0084 & 23 & 27 & 0 & $1.62 \times 10^{-6}$ \\
200 & 1.0156 & 1.0075 & 33 & 17 & 0 & $9.22 \times 10^{-9}$ \\
300 & 1.0138 & 1.0065 & 41 &  9 & 0 & $5.44 \times 10^{-10}$ \\
500 & 1.0128 & 1.0057 & 50 &  0 & 0 & $1.86 \times 10^{-10}$ \\
\midrule
\textbf{Total} & --- & --- & \textbf{162} & \textbf{88} & \textbf{0} & $< 10^{-33}$ \\
\bottomrule
\end{tabular}
\end{table}

At $n = 500$, CAD wins on \emph{every single instance}. The aggregate Wilcoxon $p$-value
is below $10^{-33}$, and CAD reduces waste by 55\% relative to the lower bound
(FFD ratio 1.0128 $\to$ CAD ratio 1.0057).

\subsection{Comprehensive Comparison}

Table~\ref{tab:comprehensive} presents results across all hard and medium instance classes.
CAD consistently outperforms all three baselines on instances where items fill a substantial
fraction of bin capacity.

\begin{table}[h]
\centering
\caption{Comprehensive comparison across instance classes ($n \geq 200$, 50 trials each).
  ``CAD$<$FFD'' and ``CAD$<$MFFD'' count instances where CAD uses strictly fewer bins.}
\label{tab:comprehensive}
\small
\begin{tabular}{llccccccc}
\toprule
\textbf{Class} & $n$ & FFD & BFD & MFFD & CAD & CAD$<$FFD & CAD$<$MFFD & $p$ \\
\midrule
\textsc{Hard-150} & 200 & 1.0156 & 1.0156 & 1.0377 & 1.0075 & 33 & 33 & $9 \times 10^{-9}$ \\
\textsc{Hard-150} & 300 & 1.0138 & 1.0138 & 1.0548 & 1.0065 & 41 & 46 & $5 \times 10^{-10}$ \\
\textsc{Hard-150} & 500 & 1.0128 & 1.0128 & 1.0645 & 1.0057 & 50 & 48 & $2 \times 10^{-10}$ \\
\textsc{Hard-200} & 300 & 1.0134 & 1.0133 & 1.0222 & 1.0081 & 33 & 42 & $3 \times 10^{-8}$ \\
\textsc{Hard-200} & 500 & 1.0114 & 1.0114 & 1.0168 & 1.0063 & 43 & 47 & $2 \times 10^{-9}$ \\
\textsc{Medium}   & 300 & 1.0117 & 1.0117 & 1.0299 & 1.0092 & 17 & 40 & $4 \times 10^{-5}$ \\
\textsc{Medium}   & 500 & 1.0071 & 1.0069 & 1.0213 & 1.0044 & 29 & 42 & $6 \times 10^{-7}$ \\
\bottomrule
\end{tabular}
\end{table}

\paragraph{Notable observations.}
\begin{enumerate}[nosep]
  \item CAD outperforms MFFD on all tested configurations. MFFD's four-category classification
        appears to degrade on hard instances, while CAD's adaptive completion captures tighter
        pairings.
  \item The improvement scales with $n$: at $n = 500$, CAD wins on 50/50 \textsc{Hard-150}
        instances.
  \item CAD never performs more than 1 bin worse than FFD on any single instance (across 1{,}400+
        instances, only 3 losses total, all by exactly 1 bin).
\end{enumerate}

\subsection{Instances Where CAD Does Not Help}

\begin{table}[h]
\centering
\caption{Instance classes where CAD provides minimal or no improvement.}
\label{tab:noimprove}
\begin{tabular}{lll}
\toprule
\textbf{Class} & \textbf{Result vs.\ FFD} & \textbf{Explanation} \\
\midrule
\textsc{Uniform} $[1,99]/100$  & Neutral & Items too small; FFD already near-optimal \\
\textsc{Large-Items} $[250,500]/1000$ & Not significant ($p = 0.17$) & 2+ items rarely fit \\
\textsc{Triplet}               & Tie & Both algorithms near-optimal on structured instances \\
\bottomrule
\end{tabular}
\end{table}

CAD's completion step is most effective when items fill approximately 13--67\% of
bin capacity, creating opportunities for 2--4 items per bin. When items are very
small (many per bin, FFD packs efficiently) or very large (at most 2 per bin,
limited completion opportunities), the advantage vanishes.

\subsection{Mechanism Analysis}

To understand \emph{why} CAD improves, we examine the waste distribution across bins
on 30 \textsc{Hard-150} instances at $n = 300$:

\begin{table}[h]
\centering
\caption{Waste analysis: FFD vs.\ CAD on \textsc{Hard-150}, $n = 300$ (30 instances).}
\label{tab:mechanism}
\begin{tabular}{lccr}
\toprule
\textbf{Metric} & \textbf{FFD} & \textbf{CAD} & \textbf{Change} \\
\midrule
Total bins          & 3{,}667 & 3{,}644 & $-23$ \\
Near-full bins (waste $< 8$) & 3{,}251 & 3{,}498 & $+247$ \\
Zero-waste bins     & 2{,}464 & 2{,}733 & $+269$ \\
Medium-waste bins (waste 8--19) & 387 & 110 & $-277$ \\
Avg.\ waste per bin & 2.55 (1.7\%) & 1.62 (1.1\%) & $-36\%$ \\
\bottomrule
\end{tabular}
\end{table}

CAD converts 277 medium-waste bins into zero-waste or near-full bins. The completion
step performs approximately 118 completions per instance (of 300 items), meaning roughly
40\% of items are placed as completions rather than through the standard BFD sequence.
This locks in good pairings before they are consumed by later bins.

\subsection{Runtime}

\begin{table}[h]
\centering
\caption{Mean wall-clock time per instance (single-threaded, Python 3.11, AMD Ryzen 7).}
\label{tab:runtime}
\begin{tabular}{rccc}
\toprule
$n$ & FFD (ms) & BFD (ms) & CAD (ms) \\
\midrule
100  & 0.10 & 0.12 & 0.80 \\
300  & 0.25 & 0.35 & 8.5 \\
500  & 0.50 & 0.70 & 35 \\
1000 & 1.20 & 2.10 & 382 \\
\bottomrule
\end{tabular}
\end{table}

CAD is approximately $10$--$35\times$ slower than FFD, but remains practical for
all instances with $n \leq 10{,}000$. A C++ implementation with a sorted
index structure could substantially reduce the constant factor.

% ============================================================================
% 5. INDEPENDENT VERIFICATION
% ============================================================================
\section{Independent Verification}

To guard against implementation artifacts, we conducted a full independent verification:
\begin{enumerate}[nosep]
  \item A second implementation of CAD, FFD, and BFD was written from scratch with
        different data structures and coding style.
  \item 200 fresh \textsc{Hard-150} instances ($n = 500$) with non-overlapping seeds were tested.
  \item Result: 160 wins, 40 ties, 0 losses ($p = 3.51 \times 10^{-31}$).
  \item A stress test of 1{,}000 random instances (varying $n$, mixed classes) found
        exactly 1 instance where CAD used 1 more bin than FFD (violation rate $\approx 0.1\%$).
\end{enumerate}

CAD is \emph{not} monotonically better than FFD---rare violations exist. However,
the expected improvement is strongly positive on all hard and medium instance classes.

% ============================================================================
% 6. RELATED WORK
% ============================================================================
\section{Related Work}

\paragraph{First-Fit Decreasing and variants.}
FFD was first analyzed by \citet{johnson1973near}, who showed an asymptotic ratio of
$\frac{11}{9}\text{OPT} + \frac{6}{9}$. This bound was later shown to be tight
\citep{dosa2007tight}. Modified FFD (MFFD) by \citet{johnson1985bin} improves the
asymptotic ratio to $\frac{71}{60}\text{OPT} + 1$ through item classification. However,
our experiments suggest MFFD can perform \emph{worse} than FFD on certain hard instances
due to its rigid classification scheme.

\paragraph{Bin completion.}
\citet{korf2002new} and \citet{fukunaga2007bin} introduced \emph{bin completion}, an
algorithm that assigns a complete set of items to each bin. This is an exact algorithm
with exponential worst-case time, whereas CAD is a greedy heuristic that performs a
single scan. The philosophical connection is that both consider filling bins rather than
placing individual items; however, CAD makes irrevocable greedy decisions and runs in
polynomial time.

\paragraph{Harmonic and asymptotic algorithms.}
Harmonic algorithms \citep{lee1985simple} and their refined variants achieve better
asymptotic ratios but are designed for the online setting. The Karmarkar--Karp
algorithm \citep{karmarkar1982efficient} achieves $\text{OPT} + O(\log^2 \text{OPT})$
bins but relies on LP rounding and is substantially more complex than CAD.

\paragraph{Metaheuristics.}
Genetic algorithms \citep{falkenauer1996hybrid}, simulated annealing, and other
metaheuristics can find near-optimal solutions but require multiple passes and
significantly more computation time. CAD is positioned as a drop-in replacement for FFD
in settings where simplicity and single-pass execution are required.

% ============================================================================
% 7. DISCUSSION
% ============================================================================
\section{Discussion}

\subsection{Limitations}

\begin{enumerate}[nosep]
  \item \textbf{No improved worst-case guarantee.} We have not proven that CAD achieves a
        better asymptotic ratio than FFD's $\frac{11}{9}$. Our evidence is purely
        empirical. An open problem is whether the completion step provably tightens the
        approximation ratio.
  \item \textbf{Not monotonically better.} CAD can use 1 more bin than FFD on $\approx 0.1\%$
        of instances. This occurs when the completion step consumes an item that would
        have been better placed later.
  \item \textbf{Instance-dependent.} The improvement vanishes on uniform small-item and
        very-large-item instances.
  \item \textbf{Threshold sensitivity.} While the default thresholds $\alpha = 0.15$,
        $\beta = 0.10$ work well empirically, they were not formally optimized.
\end{enumerate}

\subsection{Practical Applications}

CAD is most beneficial in settings where:
\begin{itemize}[nosep]
  \item Items fill 13--67\% of container capacity (2--4 items per container).
  \item The cost of an extra bin/container is significant (shipping, cloud VMs, warehouse storage).
  \item $O(n^2)$ runtime is acceptable (typical for $n < 10{,}000$).
  \item Simplicity and single-pass execution are valued over optimal but complex algorithms.
\end{itemize}

Industries where these conditions hold include logistics (container loading), cloud computing
(virtual machine bin packing), manufacturing (cutting stock), and 3D printing (nesting).

\subsection{Open Problems}

\begin{enumerate}[nosep]
  \item Does CAD achieve a provably better asymptotic approximation ratio than $\frac{11}{9}$?
  \item Can the thresholds $\alpha, \beta, \gamma$ be set adaptively based on the item size
        distribution to avoid the rare cases where CAD underperforms?
  \item Does the completion principle extend to higher-dimensional packing problems
        (2D strip packing, 3D container loading)?
\end{enumerate}

% ============================================================================
% 8. CONCLUSION
% ============================================================================
\section{Conclusion}

We presented Completion-Aware Decreasing (CAD), a simple augmentation of Best-Fit
Decreasing that searches for completing items immediately after each placement.
On hard bin packing instances, CAD achieves statistically significant improvements
over FFD, BFD, and MFFD while maintaining $O(n^2)$ time complexity. The algorithm
is easy to implement, requires no external solvers, and can serve as a drop-in
replacement for FFD in applications where moderate item-to-capacity ratios are common.

\paragraph{Reproducibility.} Complete source code (Python) and all experimental data are
available at \url{https://github.com/KathanS/cad-bin-packing}.

% ============================================================================
% REFERENCES
% ============================================================================
\bibliographystyle{plainnat}
\bibliography{references}

% ============================================================================
% APPENDIX
% ============================================================================
\appendix
\section{Full Experimental Data}
\label{app:full}

Table~\ref{tab:full} presents complete results for all tested configurations.

\begin{table}[h]
\centering
\caption{Complete results for all (class, $n$) configurations (50 trials each).}
\label{tab:full}
\small
\begin{tabular}{llccccccc}
\toprule
\textbf{Class} & $n$ & FFD & BFD & CAD & Wins & Ties & Losses & $p$-value \\
\midrule
\textsc{Hard-150} & 60  & 1.0208 & 1.0208 & 1.0085 & 15 & 35 & 0 & $1.1 \times 10^{-4}$ \\
\textsc{Hard-150} & 100 & 1.0198 & 1.0198 & 1.0084 & 23 & 27 & 0 & $1.6 \times 10^{-6}$ \\
\textsc{Hard-150} & 200 & 1.0156 & 1.0156 & 1.0075 & 33 & 17 & 0 & $9.2 \times 10^{-9}$ \\
\textsc{Hard-150} & 300 & 1.0138 & 1.0138 & 1.0065 & 41 &  9 & 0 & $5.4 \times 10^{-10}$ \\
\textsc{Hard-150} & 500 & 1.0128 & 1.0128 & 1.0057 & 50 &  0 & 0 & $1.9 \times 10^{-10}$ \\
\textsc{Hard-200} & 100 & 1.0271 & 1.0271 & 1.0217 & 14 & 34 & 2 & $2.7 \times 10^{-3}$ \\
\textsc{Hard-200} & 200 & 1.0160 & 1.0160 & 1.0113 & 19 & 31 & 0 & $2.9 \times 10^{-5}$ \\
\textsc{Hard-200} & 300 & 1.0134 & 1.0133 & 1.0081 & 33 & 17 & 0 & $3.1 \times 10^{-8}$ \\
\textsc{Hard-200} & 500 & 1.0114 & 1.0114 & 1.0063 & 43 &  7 & 0 & $2.3 \times 10^{-9}$ \\
\textsc{Medium}   & 100 & 1.0278 & 1.0278 & 1.0246 &  8 & 41 & 1 & $3.9 \times 10^{-2}$ \\
\textsc{Medium}   & 200 & 1.0170 & 1.0170 & 1.0143 & 12 & 38 & 0 & $4.9 \times 10^{-4}$ \\
\textsc{Medium}   & 300 & 1.0117 & 1.0117 & 1.0092 & 17 & 33 & 0 & $3.7 \times 10^{-5}$ \\
\textsc{Medium}   & 500 & 1.0071 & 1.0069 & 1.0044 & 29 & 20 & 1 & $6.0 \times 10^{-7}$ \\
\textsc{Large}    & 100 & 1.1070 & 1.1070 & 1.1015 & 17 & 21 & 12 & $1.7 \times 10^{-1}$ \\
\bottomrule
\end{tabular}
\end{table}

\end{document}
